\chapter{Introdução}

Este material se propõe a explorar o conceito e as propriedades de uma ferramenta bastante eficaz na resolução de problemas matemáticos das mais diversas áreas, como álgebra, combinatória e teoria dos números: \textit{polinômios simétricos}. Tal recurso se revela particularmente útil no contexto olímpico e de vestibulares militares, permitindo que se encontrem soluções simples para questões aparentemente complexas por meio da identificação de simetrias e, assim, poupando tempo precioso. No decorrer do texto, serão abordadas as \textit{Relações de Girard} e as \textit{Somas de Potências de Newton}, imprescindíveis para o entendimento do tema, e mostradas aplicações práticas dos \textit{polinômios simétricos}.

\begin{definicao}[Polinômios Simétricos]
	Um polinômio $P(x_1, x_2, \dots, x_n)$ é \textit{simétrico} se permanece invariante sob qualquer permutação, representada por $x_\sigma$, de suas variáveis:
	\begin{equation*}
		P(x_{\sigma(1)}, x_{\sigma(2)}, \dots, x_{\sigma(n)}) = P(x_1, x_2, \dots, x_n).
	\end{equation*}
\end{definicao}

\begin{teorema}[Teorema Fundamental dos Polinômios Simétricos]
Qualquer polinômio simétrico pode ser descrito em termos dos polinômios simétricos elementares. Que são:
	\begin{align*}
	    \sigma_1 &= \text{Soma das variáveis} \\
	    \sigma_2 &= \text{Soma dos produtos dois a dois das variáveis} \\
	    \sigma_3 &= \text{Soma dos produtos três a três das variáveis} \\
	             & \vdots \\
	    \sigma_n &= \text{Produto de todas as variáveis}
	\end{align*}
\end{teorema}

Pegando como exemplo $D(x, y, z) = x^2 + y^2 + z^2$, seus polinômios elementares serão:
\begin{align*}
	\sigma_1 & = x + y + z\\
	\sigma_2 & = xy + xz + yz \\
	\sigma_3 & = xyz
\end{align*}

Portanto $D(x, y, z) \text{ pode ser descrito como } \sigma_1^2 - 2\sigma_2$.

A teoria a seguir demonstra como as Relações de Girard e as Identidades de Newton operam como duas faces da mesma estrutura algébrica. Essa estrutura não apenas é elegante do ponto de vista matemático, mas também possui aplicações diretas em sistemas de equações, desigualdades e problemas de otimização.

\chapter{Relações de Girard}

As relações de Girard estabelecem uma conexão direta entre os coeficientes de um polinômio e suas raízes. No caso de polinômios simétricos, essa relação está intrinsecamente ligada ao \textit{Teorema Fundamental dos Polinômios Simétricos}.

Considere o polinômio do $2^\circ$ grau, com $a \neq 0$:
\begin{equation}
	P(x) = ax^2 + bx + c
\end{equation}

Sendo $x_1$ e $x_2$ as raízes desse polinômio, podemos escrevê-lo na forma fatorada:
\begin{align*}
	a(x - x_1)(x - x_2) & = a(x^2 - x_1x - x_2x + x_1x_2)  \\
	                    & = a(x^2 - (x_1 + x_2)x + x_1x_2) \\
	                    & = ax^2 - a(x_1 + x_2)x + ax_1x_2
\end{align*}

Definindo os polinômios elementares simétricos $\sigma_1 = x_1 + x_2$ e $\sigma_2 = x_1x_2$, observamos por comparação com a forma original que:
\begin{align*}
	b = -a\sigma_1 & \implies \sigma_1 = -\frac{b}{a} \\
	c = a\sigma_2  & \implies \sigma_2 = \frac{c}{a}
\end{align*}

Generalizando para um polinômio de grau $n$ ($n \geq 1$) com $a_n \neq 0$:
\begin{align*}
	P(x) & = a_n x^n + a_{n-1} x^{n-1} + \dots + a_1 x + a_0 \\
	     & = a_n(x - x_1)(x - x_2)\dots(x - x_n)
\end{align*}

Ao expandir o produto das raízes, obtemos a relação geral em termos dos polinômios elementares simétricos $\sigma_k$:
\begin{equation}
	P(x) = a_n \sum_{k=0}^{n} (-1)^k \sigma_k x^{n-k}
\end{equation}
onde $\sigma_0 = 1$. Portanto, para cada $k \in \{1, \dots, n\}$, temos:
\begin{equation}
	\sigma_k = (-1)^k \frac{a_{n-k}}{a_n}
\end{equation}

A elegância dessas relações torna-se ainda mais evidente na análise de somas de potências de raízes, tópico que exploraremos no capítulo seguinte.

\chapter{Soma de potências de Newton}

Em 1707, o matemático e físico Sir Isaac Newton generalizou em sua obra \textit{Arithmetica Universalis} as somas de potências, hoje conhecidas como Somas de Newton na forma:

$$S_{k} = r_{1}^{k} + r_{2}^{k} + r_{3}^{k} + \dots + r_{n}^{k}$$

Onde $r_{1}, r_{2}, \dots, r_{n}$ são raízes de $S(x)=0$. Essa soma pode ser reescrita com polinômios simétricos elementares. Tal relação é um teorema de Identidade de Newton.

Seja $k \in \mathbb{N}$, $k \ge n$, então $S(x) = x^{n} - \sigma_{1}x^{n-1} + \sigma_{2}x^{n-2} - \dots + (-1)^{n}\sigma_{n}$

Demonstração:

Usando as relações de Girard, podemos escrever a equação para uma raiz:

$$P(x_{1}) = x_{1}^{n} - \sigma_{1}x_{1}^{n-1} + \sigma_{2}x_{1}^{n-2} - \sigma_{3}x_{1}^{n-3} + \dots + (-1)^{n}\sigma_{n} = 0$$

Multiplicando por $x_{1}^{k-n}$ obtemos:

$$x_{1}^{k} - \sigma_{1}x_{1}^{k-1} + \sigma_{2}x_{1}^{k-2} - \sigma_{3}x_{1}^{k-3} + \dots + (-1)^{n}\sigma_{n}x_{1}^{k-n} = 0$$

Para cada uma das raízes montamos um sistema de equações:

$$x_{1}^{k} - \sigma_{1}x_{1}^{k-1} + \sigma_{2}x_{1}^{k-2} - \sigma_{3}x_{1}^{k-3} + \dots + (-1)^{n}\sigma_{n}x_{1}^{k-n} = 0$$

$$x_{2}^{k} - \sigma_{1}x_{2}^{k-1} + \sigma_{2}x_{2}^{k-2} - \sigma_{3}x_{2}^{k-3} + \dots + (-1)^{n}\sigma_{n}x_{2}^{k-n} = 0$$

$$\vdots$$

$$x_{n}^{k} - \sigma_{1}x_{n}^{k-1} + \sigma_{2}x_{n}^{k-2} - \sigma_{3}x_{n}^{k-3} + \dots + (-1)^{n}\sigma_{n}x_{n}^{k-n} = 0$$

Somando as equações e escrevendo em somas de Newton:

$$S_{k} - \sigma_{1}S_{k-1} + \sigma_{2}S_{k-2} - \sigma_{3}S_{k-3} + \dots + (-1)^{n}\sigma_{n}S_{k-n} = 0$$

\chapter{Aplicações}

Existem aplicações dos polinômios simétricos tanto na física quanto na criptografia, nesse texto exploraremos sua utilidade nas olimpiadas de conhecimento. A seguir temos um exeplo de quetão da OBM que podemos utilizar as propriedades desenvolvidas até aqui.

(OBM 2010) Sejam $a$, $b$ e $c$ reais tais que $a \neq b$ e
\[
a^2(b + c) = b^2(c + a) = 2010.
\]

Calcule $c^2(a + b)$
